\clearpage
\ssscp{\tsc{Havu-Dhao}}{03-03-02-HD}
Within \tsc{Sumba-Havu} there is very strong
evidence for placing Havu and Dhao together and apart
from the languages of Sumba.
This evidence comes from a process of vowel metathesis
associated with syllable-final *a or *ə.
The discussion in this section draws on
\citet{Grimes2008,GrimesRanohAplugi2008,Grimes2010a} and \citet{Blust2012b}.

We begin by discussing the more straightforward
sources of schwa in \tsc{Havu-Dhao}.
Keep in mind in all the following examples
that penultimate schwa in Havu and Dhao is syllabic,
stress bearing, and triggers phonetic
lengthening of the following consonant.

One source of schwa in \tsc{Havu-Dhao} is retention.
PMP *ə = \it{ə} in penultimate syllables,
as shown in \trf{02-09a}.

\begin{table}
	\centering\tcp{\tsc{Havu-Dhao} *ə = \it{ə}}{02-09a}
%\begin{adjustbox}{width=\textwidth}
	\st{6pt}\begin{tabular}{llllllll}\lsptoprule
gloss	&	buy	&	PP	&	egg	&	hear	&	turtle	&	bend, fold	&	fathom	\\
PMP	&	*b\tbr{ə}li	&	*\tbr{ə}(m)pu	&	*qat\tbr{ə}luR	&	*d\tbr{ə}ŋəR	&	*p\tbr{ə}ñu	&	*l\tbr{ə}ku	&	*d\tbr{ə}pa	\\ \midrule
Havu	&	{\hr}v\tbr{ə}li	&	{\hr}\tbr{ə}pu	&	\hp{*qa}d\tbr{ə}lu	&	{\hr}r\tbr{ə}ŋe	&	\hp{*p}\tbr{ə}ɲu	&	{\hr}l\tbr{ə}ku	&	{\hr}r\tbr{ə}pa	\\
Dhao	&	{\hr}h\tbr{ə}li	&	{\hr}\tbr{ə}pu	&	\tcg{(kanaɖ͡ʐu)}	&	{\hr}tad\tbr{ə}ŋi	&	\hp{*p}\tbr{ə}ɲu	&	\tcg{(ləpa)}	&	{\hr}r\tbr{ə}pa	\\
		\lspbottomrule
	\end{tabular}
	%\end{adjustbox}
\end{table}

Many penultimate schwas in Havu
and Dhao are derived from PMP *a
in proximity to high vowels *i and *u.
The most straightforward examples involve *a > \it{ə}
/{\gap}V[\tsc{+high}], including after loss of medial *q or *h.
Examples are given in \trf{02-10a}.
When an intervening consonant occurs penultimate *a = \it{a}.
Two examples are *talih > Havu \it{dari}, Dhao \it{ɖ͡ʐari} `rope, cord'
and *walu > Havu, Dhao \it{aru} `eight'.

\begin{table}
	\centering\tcp{\tsc{Havu-Dhao} *a > \it{ə} /{\gap}V\tsc{+high}}{02-10a}
	%\begin{adjustbox}{width=\textwidth}
	\st{6pt}\begin{tabular}{llllllll}\lsptoprule
*gloss	&	odour	&	person	&	\hp{*ka}far	&	leaf	&	year	&	faeces	&	water	\\	
PMP	&	*b\tbr{a}hu	&	*t\tbr{a}u	&	\hp{*ka}z\tbr{a}uq	&	*d\tbr{a}hun	&	*t\tbr{a}qun	&	*t\tbr{a}qi	&	*w\tbr{a}hiyəR	\\ \midrule	
Havu	&	{\hr}v\tbr{ə}u	&	{\hr}d\tbr{ə}u	&	\hp{*ka}ʄ\tbr{ə}u	&	{\hr}r\tbr{ə}u	&	{\hr}t\tbr{ə}u	&	{\hr}d\tbr{ə}i	&	{\hr}\tbr{ə}i	\\	
Dhao	&	{\hr}h\tbr{ə}u	&	{\hr}ɖ͡ʐ\tbr{ə}u	&	{\hr}kaʄ\tbr{ə}u	&	{\hr}r\tbr{ə}u	&	{\hr}t\tbr{ə}u	&	{\hr}ɖ͡ʐ\tbr{ə}i	&	{\hr}\tbr{ə}i	\\	
		\lspbottomrule
	\end{tabular}
	%\end{adjustbox}
\end{table}

A more complex source of schwa is from final *a or *ə
where the penultimate vowel is high, in which case vowel
metathesis also occurs;
*uCa/*uCə > \it{əCu} and *iCa/*iCə > \it{əCi}.
Examples are given in \trf{02-10b},
with Kambera and Kodi cognates to show that
this change separates \tsc{Havu-Dhao}
from \tsc{Sumba.}
%and many more Havu examples
%can be found in \citet{Blust2012a}.

\begin{table}
	\centering\tcp{\tsc{Havu-Dhao} *V\tsc{[+high]}CV\tsc{[-high]} > əC\tsc{[+high]}}{02-10b}
	%\begin{adjustbox}{width=\textwidth}
	\st{3pt}\begin{tabular}{llllllll}\lsptoprule
*gloss	&	moon									&	house										&	navel										& clay pot								&	oil										&	five									&	black	\\	
PMP			&	*b\tbb{u}l\tbr{a}n		&	*R\tbb{u}m\tbr{a}q			&	*p\tbb{u}s\tbr{ə}j			&	*k\tbb{u}d\tbr{ə}n			&	*m\tbb{i}ñ\tbr{a}k		&	*l\tbb{i}m\tbr{a}			&	*ma-q\tbb{i}t\tbr{ə}m	\\ \midrule	
Havu		&	{\hr}v\tbr{ə}r\tbb{u}	&	\hp{*R}\tbr{ə}m\tbb{u}	&	\hp{*p}\tbr{ə}h\tbb{u}	&	\hp{*k}\tbr{ə}r\tbb{u}	&	{\hr}m\tbr{ə}ɲ\tbb{i}	&	{\hr}l\tbr{ə}m\tbb{i}	&	{\hr}m\tbr{ə}d\tbb{i}	\\	
Dhao		&	{\hr}h\tbr{ə}r\tbb{u}	&	\hp{*R}\tbr{ə}m\tbb{u}	&	\hp{*p}\tbr{ə}s\tbb{u}	&	\hp{*k}\tbr{ə}r\tbb{u}	&	{\hr}m\tbr{ə}ɲ\tbb{i}	&	{\hr}l\tbr{ə}m\tbb{i}	&	{\hr}m\tbr{ə}ɖ͡ʐ\tbb{i}	\\	\midrule
Kodi		&	{\hr}wulːa						&	\hp{*R}umːa							&	‹puhu›									&	{\hr}ɣurːa							&	{\hr}miɲːa						&	{\hr}limːa						&	{\hr}mete	\\	
Kambera	&	{\hr}wulaŋ-u					&	\hp{*R}uma							&	{\hr}puhu								&	{\hr}wuruŋ-u						&	{\hr}mina							&	{\hr}lima							&	{\hr}mitiŋ-u	\\	
		\lspbottomrule
	\end{tabular}
	%\end{adjustbox}
\end{table}

In final syllables usually *ə > \it{e, a}
when the penultimate vowel is *ə or *a.
Examples are given in \trf{02-09b}.
*ə > \it{i} in Dhao *dəŋəR > \it{tadəŋi} may due to earlier *R > **y
(see discussion associated \trf{03-51}).
%XXX back reference.mention otehr sources of e

\begin{table}
	\centering\tcp{\tsc{Havu-Dhao} *ə > \it{ə, a} /{\gap}C{\#}}{02-09b}
%\begin{adjustbox}{width=\textwidth}
	\st{6pt}\begin{tabular}{lllllll}\lsptoprule
*gloss	&	bury	&	hear	&	soak	&	inside	&	six	&	sick, ill	\\
PMP	&	*tan\tbr{ə}m	&	*dəŋ\tbr{ə}R	&	*(R)əd\tbr{ə}m	&	*dal\tbr{ə}m	&	*ən\tbr{ə}m	&	*hapəj\tbr{ə}s	\\ \midrule
Havu	&	pe-dan\tbr{e}	&	{\hr}rəŋ\tbr{e}	&	\hp{(R)}\cg{(lie)}	&	{\hr}ɗar\tbr{a}	&	{\hr}ən\tbr{a}	&	\hp{*ha}pəɗ\tbr{a}	\\
Dhao	&	pa-ɖ͡ʐan\tbr{e}	&	{\hr}tadəŋ\tbr{i}	&	\hp{*(R)}eɖ͡ʐ\tbr{e}	&	{\hr}ɗar\tbr{a}	&	{\hr}ən\tbr{a}	&	\hp{*ha}pəɗ\tbr{a}, pəd\tbr{a}	\\
		\lspbottomrule
	\end{tabular}
	%\end{adjustbox}
\end{table}

Similarly, the vowel-glide sequence *-ay > \it{e} in \tsc{Havu-Dhao}.
And similarly vowel sequences *ai/*a-i > \it{e}.
Examples are given in \trf{02-09c}.
Final *-ay > \it{e} means that there is a partial merger
of *-ay and *ə in final syllables.

\begin{table}
	\centering\tcp{Sources of \it{ə} and \it{e} in Havu and Dhao}{02-09c}
%\begin{adjustbox}{width=\textwidth}
	\st{6pt}\begin{tabular}{llllllllll}\lsptoprule
gloss	& ascend				&rice							&liver						&die							&man, male				& red-brown\\
PMP		&*sak\tbr{ay}		&*paj\tbr{ay}			&*qat\tbr{ay}			&*mat\tbr{ay}			&*maRuqan\tbr{ay}	&*m\tbr{a-i}Raq\\ \midrule
Havu	&{\hr}haʔ\tbr{e}&\hp{*p}ar\tbr{e}	&\hp{*q}ad\tbr{e}	&{\hr}mad\tbr{e}	&{\hr}mon\tbr{e}	&{\hr}m\tbr{e}a\\
Dhao	&{\hr}ʧaʔ\tbr{e}&\hp{*p}ar\tbr{e}	&\hp{*q}aɖ͡ʐ\tbr{e}&{\hr}maɖ͡ʐ\tbr{e}	&{\hr}mon\tbr{e}	&{\hr}m\tbr{e}a\\
		\lspbottomrule
	\end{tabular}
%	\end{adjustbox}
\end{table}
